\documentclass[10pt,a4paper]{amsart}
\usepackage[margin=19mm]{geometry}
\usepackage{amsmath,amssymb,mathtools}
\usepackage[scr=rsfs,cal=euler]{mathalfa}
\usepackage{xfrac}
\usepackage[unicode,hidelinks,bookmarks=false]{hyperref}
\newcommand{\ie}{i.e.\ }
\newcommand{\cf}{cf.\ }
\newcommand{\resp}{respectively\ }
\newcommand{\Subsection}{\subsection}
\providecommand{\nobreakdash}{\nobreak\hbox{-}}
\setlength{\emergencystretch}{2em}
\multlinegap0pt
\usepackage{mathtools}
\usepackage{amssymb}
\usepackage{xcolor}
\usepackage{braket}
\usepackage{dsfont}
\usepackage{algorithm}\usepackage{algpseudocode}
\usepackage[shortlabels]{enumitem}
\usepackage{float}
\usepackage[normalem]{ulem}
\usepackage{tikz}
\newtheorem{lemma}{Lemma}
\newtheorem{theorem}{Theorem}
\usepackage{cleveref}
\crefname{lemma}{Lemma}{Lemmas}
\crefname{theorem}{Theorem}{Theorems}
\newcommand{\KK}{\mathbb{K}}
\newcommand{\PP}{\mathbb{P}}
\DeclareMathOperator{\Res}{Res}
\newcommand{\cZ}{\mathcal{Z}}
\title{Small Resultant Systems via Linear Combinations}
\author{M. Levent Do\u{g}an, Elias Tsigaridas, Zafeirakis Zafeirakopoulos}
\date{Source: arXiv:2607.29308v1 (CC BY 4.0)}
\begin{document}
\maketitle

\noindent\emph{Source and licence.} Small Resultant Systems via Linear Combinations. Version arXiv:2607.29308v1 (CC BY 4.0), distributed by arXiv under the Creative Commons Attribution 4.0 International licence, http://creativecommons.org/licenses/by/4.0/, as recorded in the arXiv licence metadata for this exact version and shown on the abstract page https://arxiv.org/abs/2607.29308v1. Full licence notice: \texttt{licenses/arxiv-2607.29308-CC-BY-4.0.txt}.\par\smallskip

\noindent\textit{Editorial scope.} Counted unit: the article's theorem thm:explicit-multivariate (source0.tex lines 1356--1368). The article's complete original proof of it is delivered with it (source0.tex lines 1369--1393, one \texttt{proof} environment of 25 lines). No mathematics is written, repaired or re-derived: the statement and the proof are the article's own lines, in the article's own notation, and no retained line is substituted or re-flowed.  Retained uncounted as the source-local context the proof needs: lines 1222--1234; lines 1322--1339.  Nothing was omitted from any retained span, so the package carries no omission record.  No retained line contains a citation, so the wrapper carries no bibliography and no bibliography entry.  No retained line contains a figure environment, an image-inclusion command or drawing code, so no image is delivered and the package declares no asset dependency.  Editorial formatting only, in addition: an AMS \texttt{amsart} preamble that restates the article's own declarations and macros used by the retained lines, namely the environments \texttt{lemma}, \texttt{theorem} on separate counters, together with the standard packages those lines need.  The retained environments print in this document as Lemma 1, Theorem 1 in order of appearance, whereas the article prints the same items with the numbering of its own full document.  The complete native source of the article as distributed in the arXiv e-print for this exact version is bundled unchanged as \texttt{sources/206470/source0.tex}.
\begin{lemma}
\label{lem:vandermonde}
Suppose $m, s>0$ and $v_1,\dots,v_s\in\KK^m$ are $s$ vectors.
If $\lambda_1,\dots,\lambda_{s}\in\KK$ are $s$ distinct scalars, then the linear span of $v_1,\dots,v_s$ coincides with the linear span of the vectors
$w_1, \dots, w_s \in \KK^m$, where
 \[
\begin{split}
w_1 &= v_1 + \lambda_1 v_2 + \dots + \lambda_1^{s-1} v_s\\
& \qquad\vdots \\
w_s &= v_1 + \lambda_s v_2 + \dots + \lambda_s^{s-1} v_s.
\end{split}
\]
\end{lemma}

\begin{lemma}
\label{prop:dimension-drop}
    Suppose $C\subset\PP^{n-1}$ is pure $m$-dimensional of degree~$D$.
    Let \[
    f^\lambda \coloneqq f_1 + \lambda f_2 + \dots + \lambda^{q-1} f_q
    \] for $q$ forms $f_1,f_2,\dots,f_q$ and $\lambda\in\KK$.
    Suppose:
    \begin{enumerate}
        \item $\dim\left(C\cap\cZ_{\PP}(f_1,f_2,\dots,f_q)\right)<m$;
        \item $C\subset\cZ_{\PP}(f^{\tau_1},f^{\tau_2},\dots,f^{\tau_r})$ for distinct scalars $\tau_1,\dots,\tau_r$;
        \item $S\subset\KK$ is disjoint from $\{\tau_1,\dots,\tau_r\}$ and satisfies \[
        |S|\geq (q-r-1)D+1.
        \]
    \end{enumerate}
    Then some $\lambda\in S$ satisfies \[
    \dim\left(\, C\cap\cZ_{\PP}(f^\lambda) \,\right) = m-1.
    \]
\end{lemma}

\begin{theorem}
\label{thm:explicit-multivariate}
Suppose $S_1,S_2,\dots,S_{n-1}\subset\KK$ are pairwise disjoint sets of scalars of cardinalities \[
\forall \ell=1,2,\dots,n-1, \; |S_{\ell}| = (s-\ell-1)d^{\ell}+1.
\] Then the following collection of polynomials is a punctured resultant system: \begin{equation}
\label{eq:explicit-system-multi}
\Res(f^{\lambda_1}, f^{\lambda_2},\dots, f^{\lambda_{n-1}}, f_s),\quad (\lambda_1,\dots,\lambda_{n-1})\in S_1 \times S_2\times\dots\times S_{n-1},
\end{equation} where we denote $f^{\lambda_i}\coloneqq f_1 + \lambda_i f_2 +\dots + \lambda_i^{s-2} f_{s-1}$.

In particular, there exists a punctured resultant system of cardinality \[
|S_1\times\dots\times S_{n-1}| = \prod_{\ell=1}^{n-1} (s-\ell-1) d^{\ell} + 1.
\] 
\end{theorem}

\begin{proof}
We need to prove that if every resultant in \eqref{eq:explicit-system-multi} vanishes, then $f_1,\dots,f_s$ have a common non-zero solution.
Since $f_s\neq 0$, the hypersurface $C_0\coloneqq \cZ_{\PP}(f_s)\subset\PP^{n-1}$ has dimension $n-2$ and its degree is bounded by $d$. 
If $\cZ(f_1,\dots,f_{s-1})\cap C_0\neq\varnothing$, then we are done.
Otherwise, an application of \cref{prop:dimension-drop} to $C_0$ and to the polynomials $f_1,\dots,f_{s-1}$ with $q=s-1$ and $r=0$ implies that there exists $\lambda^{(1)}\in S_1$ such that \[
\dim\left(\cZ(f_s,f^{\lambda^{(1)}})\right) = n-3.
    \] 
    Hence, the intersection $C_1\coloneqq\cZ_{\PP}(f_s,f^{\lambda^{(1)}})$ is a pure $(n-3)$-dimensional projective algebraic variety and its degree is bounded by $d^2$ by B\'{e}zout's theorem.
    If $C_1$ intersects $\cZ_{\PP}(f_1,\dots,f_{s-1})$, then we are done. 
    Otherwise, another application of \cref{prop:dimension-drop} to $C_1$ with $q=s-1$ and $r=1$ yields a $\lambda^{(2)}\in S_2$ such that \[
    C_2\coloneqq \cZ_{\PP}(f_s,f^{\lambda^{(1)}},f^{\lambda^{(2)}}),\quad \dim C_2 = n-4.
    \] By iterating this process, we can find points $\lambda^{(i)}\in S_i$ such that \[
    C_{n-2}\coloneqq \cZ_{\PP}(f_s,f^{\lambda^{(1)}},f^{\lambda^{(2)}},\dots,f^{\lambda^{(n-2)}}),\quad \dim C_{n-2} = 0.
    \] By B\'{e}zout's theorem, the cardinality of $C_{n-2}$ is at most $\leq d^{n-1}$.

    By the assumption, the resultant \[
    \Res(f^{\mu},f^{\lambda^{(1)}},f^{\lambda^{(2)}},\dots,f^{\lambda^{(n-2)}},f_s) = 0
    \] vanishes for every $\mu\in S_{n-1}$.
    Hence, for every $\mu\in S_{n-1}$ there exists a point $p\in C_{n-2}$ such that $f^{\mu}$ vanishes on $p$.
    By pigeonhole principle, there exist $\mu_1,\dots,\mu_{s-n+1}\in S_{n-1}$ such that $f^{\mu_1},\dots,f^{\mu_{s-n+1}}$ vanish on the same point $p\in C_{n-2}$. 
    By \cref{lem:vandermonde},
    the polynomials $f^{\mu_1},\dots,f^{\mu_{s-n+1}}, f^{\lambda^{(1)}},\dots,f^{\lambda^{(n-2)}}$ have the same linear span as $f_1,\dots,f_{s-1}$.
    Hence, $p$ is in the common zero locus of $f_1,\dots,f_{s-1},f_s$.
    This finishes the proof.
\end{proof}

\end{document}
